% STANDALONE, self-contained version of resonance_sweep_table.tex.
%
% Unlike every other .tex file in analysis/ (which are fragments meant to
% be pasted into the thesis's main.tex and rely on that document's existing
% preamble), this one is a complete, compilable document on its own -
% \documentclass, full preamble, \begin{document}...\end{document}. Open
% and compile this single file directly (e.g. as its own standalone
% Overleaf project) to preview/verify just this table, with zero
% dependency on the thesis document or any other table file in this repo.
%
% To actually use this table IN the thesis, copy the content between
% \begin{document} and \end{document} into main.tex where the table
% should appear (matching how every other analysis/*_table.tex file is
% used) - don't \input{} this whole file there, since it duplicates
% preamble packages main.tex already loads.
%
% See resonance_sweep_table.tex for the full methodology/data-provenance
% header comment this table is built from (not repeated here to avoid the
% two files drifting out of sync in the parts that do matter - the actual
% table content below IS kept identical between the two).

\documentclass[11pt]{article}
\usepackage[margin=1in]{geometry}
\usepackage{booktabs}
\usepackage{siunitx}
\usepackage{xcolor}
\usepackage{colortbl}

\begin{document}

\definecolor{paramgrey}{gray}{0.85}

\begin{table}[h]
\centering
\tiny
\caption{Resonance duck mode parameter sweep (100-point Latin Hypercube,
Construction Scene, Dialogue priority): five representative points
isolating bandwidth, not peak count, as the parameter that separates a
strong result from a weak one. Gains are relative to the unprocessed
baseline. Best point found here, ranked by STOI gain (margin +3.67dB,
STOI +0.070, loudness reduction +10.6pp), falls meaningfully short of
Advanced mode's own 150-point-sweep ceiling on the same scene (margin
+6.96dB, STOI +0.104) -- loudness reduction is the one metric where
Resonance is competitive (this sweep's own loudness-reduction ceiling,
+11.5pp, essentially matches Advanced's +11.4pp).}
\label{tab:resonance-sweep}
\begin{tabular}{@{}l r r r r r r r r r r@{}}
\toprule
\cellcolor{paramgrey}\textbf{Point} & \cellcolor{paramgrey}\textbf{Thr.} & \cellcolor{paramgrey}\textbf{Ratio} & \cellcolor{paramgrey}\textbf{Knee} & \cellcolor{paramgrey}\textbf{Atk.} & \cellcolor{paramgrey}\textbf{Rel.} & \cellcolor{paramgrey}\textbf{Peaks} & \cellcolor{paramgrey}\textbf{BW (oct)} & \textbf{Margin Gain} & \textbf{STOI Gain} & \textbf{Loud.\ Red.} \\
 & \cellcolor{paramgrey}(dB) & \cellcolor{paramgrey}{} & \cellcolor{paramgrey}(dB) & \cellcolor{paramgrey}(ms) & \cellcolor{paramgrey}(ms) & \cellcolor{paramgrey}{} & \cellcolor{paramgrey}{} & (dB) & & (pp) \\
\midrule
Best overall (3 peaks)   & -50.6 & 2.4 & 11.4 & 46.0 & 319 & 3 & 3.98 & \cellcolor{green!70}+3.67 & \cellcolor{yellow!70}+0.070 & \cellcolor{orange!66}+10.6 \\
Best 1-peak              & -56.0 & 6.1 & 5.0  & 13.9 & 230 & 1 & 3.56 & \cellcolor{green!50}+2.36 & \cellcolor{yellow!66}+0.065 & \cellcolor{orange!70}+11.5 \\
Best 2-peak              & -52.4 & 8.8 & 23.6 & 16.8 & 381 & 2 & 2.28 & \cellcolor{green!41}+1.76 & \cellcolor{yellow!46}+0.040 & \cellcolor{orange!54}+8.1  \\
Narrow bandwidth (contrast) & -54.7 & 9.3 & 6.5  & 24.3 & 155 & 4 & 0.12 & \cellcolor{green!16}+0.08 & \cellcolor{yellow!20}+0.006 & \cellcolor{orange!20}+1.0  \\
Weak threshold (floor)   & -18.2 & 6.0 & 13.7 & 29.3 & 174 & 2 & 2.73 & \cellcolor{green!15}+0.00 & \cellcolor{yellow!15}+0.000 & \cellcolor{orange!15}+0.0  \\
\bottomrule
\end{tabular}
\par\vspace{4pt}
\raggedright\footnotesize Note the "Narrow bandwidth" row uses a threshold
about as aggressive as the top three rows (\SI{-54.7}{\decibel}) and a
moderate peak count (4) -- what collapses its result to almost nothing is
specifically its narrow \SI{0.12}{octave} bandwidth, not its threshold or
peak count. Across the full 100-point sweep, bandwidth is the
second-strongest predictor of STOI gain after threshold (standardized
regression coefficient $+0.505$, vs.\ threshold's $-0.775$, $R^2=0.813$),
while peak count's independent contribution is negligible ($+0.061$) --
1--2 peaks and 3--6 peaks perform statistically indistinguishably once
bandwidth is already wide. Margin gain is a weaker signal for Resonance
mode specifically than for Basic/Advanced: it is measured on Dialogue's
fixed nominal band, but Resonance ducks wherever the key's spectral peaks
currently are, not a fixed band -- see
\texttt{analysis/resonance\_sweep\_findings.txt} for the full caveat.
\end{table}

\end{document}
